\begin{figure}[htbp]
\centering
\begin{tikzpicture}[x=1cm,y=1cm]
\node[dissertation block,text width=3.65cm] (s1) at (-4.70,1.75)
  {1. Попереднє оброблення\\$S_0\rightarrow S_1$\\$X_{mm}\rightarrow X_c$};
\node[dissertation block,text width=3.65cm] (s2) at (0,1.75)
  {2. Формування ознак\\$S_1\rightarrow S_2$\\$H,J_b$};
\node[dissertation block,text width=3.65cm] (s3) at (4.70,1.75)
  {3. Відбір та інтеграція\\$S_2\rightarrow S_3$\\$X_b,Z$};
\node[dissertation block,text width=4.15cm] (s4) at (2.40,-1.40)
  {4. Прийняття рішення та визначення припису\\$S_3\rightarrow S_4$\\$S_4=\langle S_A,D_A\rangle$};
\node[dissertation block,text width=4.15cm] (s5) at (-2.40,-1.40)
  {5. Реєстрація результатів\\$S_4\rightarrow S_5$\\$S_5=F_R(S_4,C_R)$};

\draw[dissertation arrow] (s1.east) -- (s2.west);
\draw[dissertation arrow] (s2.east) -- (s3.west);
\draw[dissertation arrow] (s3.south) -- (s4.north);
\draw[dissertation arrow] (s4.west) -- (s5.east);
\end{tikzpicture}
\caption{П’ять етапів інформаційної технології виявлення ознак аномального вмісту в мультимедійних даних вебсередовища}
\label{fig:information_technology}
\end{figure}
